% LaTeX Template for AI Problem Analysis Output (v2.0)
% Strategic Analysis of AMC12B 2023 Problem 25
% IMPORTANT: Use LaTeX commands for all mathematical symbols, NOT unicode characters

\documentclass[12pt]{article}
\usepackage[utf8]{inputenc}
\usepackage{amsmath, amssymb, amsthm}
\usepackage{geometry}
\usepackage{xcolor}
\usepackage[most]{tcolorbox}
\usepackage{enumitem}
\usepackage{fancyhdr}

% Page setup
\geometry{margin=1in}
\setlength{\headheight}{14.5pt}
\pagestyle{fancy}
\fancyhf{}
\rhead{AMC12B 2023 Problem 25 Analysis}
\lhead{\today}

% Color definitions
\definecolor{problemconfig}{RGB}{230, 242, 255}    % Light blue
\definecolor{strategic}{RGB}{255, 230, 242}       % Light pink
\definecolor{tactical}{RGB}{242, 255, 230}        % Light green
\definecolor{techniques}{RGB}{255, 242, 230}      % Light yellow
\definecolor{verification}{RGB}{255, 230, 230}    % Light red
\definecolor{solution}{RGB}{230, 255, 255}        % Light cyan
\definecolor{competition}{RGB}{242, 242, 255}     % Light purple
\definecolor{proofwriting}{RGB}{240, 230, 255}    % Light lavender
\definecolor{gold}{RGB}{218, 165, 32}

% Box styles
\newtcolorbox{problemconfigbox}{
    colback=problemconfig,
    colframe=blue,
    title=Problem Configuration Analysis,
    fonttitle=\bfseries,
    arc=3pt,
    boxrule=1pt,
    breakable
}

\newtcolorbox{strategicbox}{
    colback=strategic,
    colframe=purple,
    title=Strategic Framework Application,
    fonttitle=\bfseries,
    arc=3pt,
    boxrule=1pt,
    breakable
}

\newtcolorbox{tacticalbox}{
    colback=tactical,
    colframe=green,
    title=Tactical Methodology Selection,
    fonttitle=\bfseries,
    arc=3pt,
    boxrule=1pt,
    breakable
}

\newtcolorbox{techniquesbox}{
    colback=techniques,
    colframe=orange,
    title=Conversion Techniques Application,
    fonttitle=\bfseries,
    arc=3pt,
    boxrule=1pt,
    breakable
}

\newtcolorbox{verificationbox}{
    colback=verification,
    colframe=red,
    title=Verification Strategy Design,
    fonttitle=\bfseries,
    arc=3pt,
    boxrule=1pt,
    breakable
}

\newtcolorbox{solutionbox}{
    colback=solution,
    colframe=cyan,
    title=Solution Roadmap,
    fonttitle=\bfseries,
    arc=3pt,
    boxrule=1pt,
    breakable
}

\newtcolorbox{competitionbox}{
    colback=competition,
    colframe=purple,
    title=Competition Strategy Integration,
    fonttitle=\bfseries,
    arc=3pt,
    boxrule=1pt,
    breakable
}

\newtcolorbox{proofwritingbox}{
    colback=proofwriting,
    colframe=purple!70,
    title=Proof-Writing Strategy,
    fonttitle=\bfseries,
    arc=3pt,
    boxrule=1pt,
    breakable
}

\newtcolorbox{keyinsightbox}{
    colback=yellow!20,
    colframe=gold,
    title=Key Insight,
    fonttitle=\bfseries,
    arc=3pt,
    boxrule=2pt,
    leftrule=5pt,
    breakable
}

\newtcolorbox{warningbox}{
    colback=red!10,
    colframe=red!50,
    title=Warning: Common Pitfall,
    fonttitle=\bfseries,
    arc=3pt,
    boxrule=2pt,
    leftrule=5pt,
    breakable
}

\newtcolorbox{tipbox}{
    colback=green!10,
    colframe=green!50,
    title=Pro Tip,
    fonttitle=\bfseries,
    arc=3pt,
    boxrule=2pt,
    leftrule=5pt,
    breakable
}

\begin{document}

\title{AMC12B 2023 Problem 25\\Strategic Analysis}
\author{Competition Math Framework v2.1}
\date{\today}
\maketitle

\section*{Problem Statement}

A regular pentagon with area $\sqrt{5}+1$ is printed on paper and cut out. The five vertices of the pentagon are folded into the center of the pentagon, creating a smaller pentagon. What is the area of the new pentagon?

$$\textbf{(A)}~4-\sqrt{5}\qquad\textbf{(B)}~\sqrt{5}-1\qquad\textbf{(C)}~8-3\sqrt{5}\qquad\textbf{(D)}~\frac{\sqrt{5}+1}{2}\qquad\textbf{(E)}~\frac{2+\sqrt{5}}{3}$$

\begin{problemconfigbox}
\subsection*{A. Problem Classification}
\begin{itemize}[leftmargin=*]
    \item \textbf{Problem Type}: To Find (Area of folded pentagon)
    \item \textbf{Competition Context}: AMC 12B 2023, Problem 25
    \item \textbf{Difficulty Band}: Late-band (P25) - Extremely challenging (final problem)
    \item \textbf{Mathematical Domain}: Geometry (Regular Pentagon, Paper Folding, Similar Figures, Golden Ratio)
    \item \textbf{Estimated Time}: 8-12 minutes (requires deep understanding of pentagon geometry)
\end{itemize}

\subsection*{B. Problem Structure Identification}
\begin{itemize}[leftmargin=*]
    \item \textbf{Given Information}: 
    \begin{itemize}
        \item Regular pentagon with area $\sqrt{5}+1$
        \item Five vertices folded into center
        \item Creates a smaller pentagon
        \item Multiple choice with specific algebraic answers involving $\sqrt{5}$
    \end{itemize}
    \item \textbf{Constraints}: 
    \begin{itemize}
        \item Pentagon is regular (all sides and angles equal)
        \item Folding operation: vertices meet at center
        \item Resulting shape is also a pentagon
        \item Paper folding creates perpendicular bisectors
    \end{itemize}
    \item \textbf{Objective}: 
    \begin{itemize}
        \item Find area of smaller pentagon after folding
        \item Express in form involving $\sqrt{5}$ (golden ratio appears in pentagons)
    \end{itemize}
    \item \textbf{Key Observations}: 
    \begin{itemize}
        \item Regular pentagons involve the golden ratio $\phi = \frac{1+\sqrt{5}}{2}$
        \item Area $\sqrt{5}+1$ suggests golden ratio relationship
        \item Paper folding: fold line is perpendicular bisector of segment from vertex to center
        \item Smaller pentagon likely similar to original by symmetry
        \item Answer choices all involve $\sqrt{5}$, confirming golden ratio involvement
    \end{itemize}
\end{itemize}

\subsection*{C. Strategic Orientation}
\begin{itemize}[leftmargin=*]
    \item \textbf{Problem Category}: Computational with geometric insight
    \item \textbf{Solution Approach}: 
    \begin{itemize}
        \item Understand paper folding geometry
        \item Establish similarity between pentagons
        \item Find scaling ratio using perpendicular bisector property
        \item Use similarity to find area ratio
        \item Apply golden ratio properties of regular pentagons
    \end{itemize}
    \item \textbf{Cross-Domain Potential}: 
    \begin{itemize}
        \item Geometry: Pentagon properties, similar figures
        \item Algebra: Golden ratio calculations
        \item Trigonometry: Special angles (36°, 54°, 72°, 108°)
    \end{itemize}
\end{itemize}
\end{problemconfigbox}

\begin{strategicbox}
\subsection*{A. Psychological Strategy}
\begin{itemize}[leftmargin=*]
    \item \textbf{Confidence Calibration}: This is P25 - expect significant complexity
    \item \textbf{First Impressions}: 
    \begin{itemize}
        \item Recognize golden ratio involvement from $\sqrt{5}$ in problem and answers
        \item Paper folding suggests perpendicular bisector property
        \item Regular pentagon → similar smaller pentagon by symmetry
    \end{itemize}
    \item \textbf{Emotional Management}: 
    \begin{itemize}
        \item Don't panic if golden ratio formulas aren't memorized
        \item Can derive key relationships from scratch if needed
        \item Focus on the folding geometry first
    \end{itemize}
\end{itemize}

\subsection*{B. Investigation Strategy}
\begin{itemize}[leftmargin=*]
    \item \textbf{Initial Exploration}: 
    \begin{itemize}
        \item Draw a regular pentagon with center O
        \item Mark one vertex A and understand: when A is folded to O, the crease is the perpendicular bisector of segment AO
        \item The crease passes through the midpoint M of AO
        \item By symmetry, all five creases form the smaller pentagon
    \end{itemize}
    \item \textbf{Pattern Recognition}: 
    \begin{itemize}
        \item Regular pentagon has circumradius $r_b$ (big) and apothem $a_b$
        \item Smaller pentagon has circumradius $r_s$ and apothem $a_s$
        \item Key insight: The fold line passes through midpoint of radius
        \item Therefore: $a_s = \frac{r_b}{2}$ (apothem of small = half radius of big)
    \end{itemize}
    \item \textbf{Relationship Discovery}: 
    \begin{itemize}
        \item For regular pentagon: $\cos(36^\circ) = \frac{\text{apothem}}{\text{circumradius}}$
        \item Golden ratio appears: $\cos(36^\circ) = \frac{\phi}{2} = \frac{1+\sqrt{5}}{4}$
        \item Can relate $r_s$ to $r_b$ using this property
    \end{itemize}
\end{itemize}

\subsection*{C. Argument Construction Strategy}
\begin{itemize}[leftmargin=*]
    \item \textbf{Approach}: Establish similarity ratio via folding geometry
    \item \textbf{Key Steps}:
    \begin{enumerate}
        \item Identify that $a_s = \frac{r_b}{2}$ from paper folding
        \item Use $\cos(36^\circ) = \frac{a_s}{r_s} = \frac{1+\sqrt{5}}{4}$ to find $r_s$
        \item Calculate scaling ratio $\frac{r_s}{r_b}$
        \item Area ratio = (scaling ratio)$^2$
        \item Apply to given area
    \end{enumerate}
    \item \textbf{Verification Points}: 
    \begin{itemize}
        \item Check that answer has correct form with $\sqrt{5}$
        \item Verify scaling ratio is less than 1 (smaller pentagon)
        \item Confirm answer matches one of the choices
    \end{itemize}
\end{itemize}

\subsection*{D. Late-Band Playbook (P25 Specific)}
\begin{itemize}[leftmargin=*]
    \item \textbf{P25 Characteristics}: 
    \begin{itemize}
        \item Final problem - requires deepest mathematical insight
        \item Often involves elegant special properties
        \item Golden ratio in pentagons is classic advanced geometry
        \item Paper folding adds layer of spatial reasoning
    \end{itemize}
    \item \textbf{Resource Allocation}: 
    \begin{itemize}
        \item Only attempt if confident with pentagon geometry
        \item Budget 8-12 minutes minimum
        \item Must know or be able to derive $\cos(36^\circ) = \frac{1+\sqrt{5}}{4}$
        \item Strong algebra skills needed for golden ratio manipulation
    \end{itemize}
    \item \textbf{When to Skip}: 
    \begin{itemize}
        \item If unfamiliar with pentagon special angles
        \item If golden ratio relationships are completely foreign
        \item If behind on time (other problems more valuable)
        \item If can't visualize the folding operation clearly
    \end{itemize}
\end{itemize}
\end{strategicbox}

\begin{tacticalbox}
\subsection*{A. Core Mathematical Tactics}
\begin{itemize}[leftmargin=*]
    \item \textbf{Symmetry Exploitation}: 
    \begin{itemize}
        \item Regular pentagon has 5-fold rotational symmetry
        \item Folding operation preserves this symmetry
        \item Therefore smaller shape is also regular pentagon
        \item Only need to analyze one vertex/fold
    \end{itemize}
    \item \textbf{Similar Figures}: 
    \begin{itemize}
        \item Both pentagons are regular → similar by definition
        \item Area ratio = (linear ratio)$^2$
        \item Need to find ratio of corresponding lengths
        \item Use circumradii: $\frac{A_s}{A_b} = \left(\frac{r_s}{r_b}\right)^2$
    \end{itemize}
    \item \textbf{Perpendicular Bisector Property}: 
    \begin{itemize}
        \item Paper folding creates perpendicular bisector
        \item When vertex A folds to center O, fold line $\perp$ AO
        \item Fold line passes through midpoint M of AO
        \item This is the key geometric constraint
    \end{itemize}
\end{itemize}

\subsection*{B. Domain-Specific Tactics}
\begin{itemize}[leftmargin=*]
    \item \textbf{Geometry - Regular Pentagon Properties}:
    \begin{itemize}
        \item Internal angle: $108^\circ$
        \item Central angle: $72^\circ$
        \item Relationship: $\cos(36^\circ) = \frac{\text{apothem}}{r}$
        \item Golden ratio appears: $\phi = \frac{1+\sqrt{5}}{2}$, $\cos(36^\circ) = \frac{\phi}{2}$
    \end{itemize}
    \item \textbf{Golden Ratio Algebra}:
    \begin{itemize}
        \item $\phi = \frac{1+\sqrt{5}}{2}$, $\phi^2 = \phi + 1$
        \item $\frac{1}{\phi} = \phi - 1 = \frac{\sqrt{5}-1}{2}$
        \item Rationalizing: $\frac{a}{\sqrt{5}+1} = \frac{a(\sqrt{5}-1)}{4}$
    \end{itemize}
    \item \textbf{Coordinate/Trigonometric Approach}:
    \begin{itemize}
        \item Can place pentagon in coordinate system
        \item Use special angles: $36^\circ, 54^\circ, 72^\circ, 108^\circ$
        \item $\sin(54^\circ) = \cos(36^\circ) = \frac{1+\sqrt{5}}{4}$
    \end{itemize}
\end{itemize}

\subsection*{C. Conversion Technique Selection}
\begin{itemize}[leftmargin=*]
    \item \textbf{High-Probability Techniques}:
    \begin{itemize}
        \item Similarity ratio method (most direct)
        \item Use folding property to establish $a_s = \frac{r_b}{2}$
        \item Apply $\cos(36^\circ)$ relationship
    \end{itemize}
    \item \textbf{Medium-Probability Techniques}:
    \begin{itemize}
        \item Coordinate geometry with explicit calculations
        \item Area subtraction (original minus folded triangles)
        \item Direct calculation using side lengths
    \end{itemize}
    \item \textbf{Answer Choice Testing}:
    \begin{itemize}
        \item Less effective here - need the geometric insight
        \item Could verify final answer makes sense
    \end{itemize}
\end{itemize}
\end{tacticalbox}

\begin{techniquesbox}
\subsection*{A. Primary Technique: Similarity Ratio via Folding Geometry}
\begin{itemize}[leftmargin=*]
    \item \textbf{Technique Name}: Similar Figures with Perpendicular Bisector Constraint
    \item \textbf{When to Apply}: 
    \begin{itemize}
        \item Regular polygons with folding operations
        \item Similarity established by symmetry
        \item Need to find scaling ratio
    \end{itemize}
    \item \textbf{Application - Complete Solution}: 
    
    \textbf{Step 1: Understand the Folding Operation}
    
    Let the original pentagon be $ABCDE$ centered at $O$, with circumradius $r_b$.
    
    When vertex $A$ is folded onto center $O$:
    \begin{itemize}
        \item The fold line is the perpendicular bisector of segment $AO$
        \item Let $M$ be the midpoint of $AO$, so $AM = MO = \frac{r_b}{2}$
        \item The fold line passes through $M$ and is perpendicular to $AO$
    \end{itemize}
    
    By symmetry, all five fold lines form a regular pentagon (the smaller pentagon).
    
    \textbf{Step 2: Identify the Apothem Relationship}
    
    The key observation: The fold line (edge of small pentagon) passes through $M$, the midpoint of the radius $OA$.
    
    The distance from $O$ to the fold line is the apothem $a_s$ of the small pentagon.
    
    Since the fold line passes through $M$ (midpoint of $OA$) and is perpendicular to $OA$:
    \begin{align*}
    a_s = OM = \frac{OA}{2} = \frac{r_b}{2}
    \end{align*}
    
    \textbf{Step 3: Use Pentagon Properties}
    
    For a regular pentagon with circumradius $r$ and apothem $a$:
    \begin{align*}
    \cos(36^\circ) = \frac{a}{r}
    \end{align*}
    
    The value $\cos(36^\circ) = \frac{1+\sqrt{5}}{4} = \frac{\phi}{2}$ where $\phi = \frac{1+\sqrt{5}}{2}$ is the golden ratio.
    
    Applying this to the small pentagon:
    \begin{align*}
    \cos(36^\circ) &= \frac{a_s}{r_s}\\
    \frac{1+\sqrt{5}}{4} &= \frac{a_s}{r_s}\\
    r_s &= \frac{4a_s}{1+\sqrt{5}}
    \end{align*}
    
    \textbf{Step 4: Substitute $a_s = \frac{r_b}{2}$}
    
    \begin{align*}
    r_s &= \frac{4 \cdot \frac{r_b}{2}}{1+\sqrt{5}}\\
    &= \frac{2r_b}{1+\sqrt{5}}
    \end{align*}
    
    \textbf{Step 5: Calculate the Scaling Ratio}
    
    \begin{align*}
    \frac{r_s}{r_b} &= \frac{2}{1+\sqrt{5}}
    \end{align*}
    
    Rationalize the denominator:
    \begin{align*}
    \frac{r_s}{r_b} &= \frac{2}{1+\sqrt{5}} \cdot \frac{1-\sqrt{5}}{1-\sqrt{5}}\\
    &= \frac{2(1-\sqrt{5})}{1-5}\\
    &= \frac{2(1-\sqrt{5})}{-4}\\
    &= \frac{\sqrt{5}-1}{2}
    \end{align*}
    
    Note: This is $\frac{1}{\phi}$ where $\phi$ is the golden ratio!
    
    \textbf{Step 6: Calculate Area Ratio}
    
    Since the pentagons are similar:
    \begin{align*}
    \frac{A_s}{A_b} &= \left(\frac{r_s}{r_b}\right)^2\\
    &= \left(\frac{\sqrt{5}-1}{2}\right)^2\\
    &= \frac{(\sqrt{5}-1)^2}{4}\\
    &= \frac{5 - 2\sqrt{5} + 1}{4}\\
    &= \frac{6 - 2\sqrt{5}}{4}\\
    &= \frac{3 - \sqrt{5}}{2}
    \end{align*}
    
    \textbf{Step 7: Apply to Given Area}
    
    Given $A_b = \sqrt{5} + 1$:
    \begin{align*}
    A_s &= A_b \cdot \frac{3-\sqrt{5}}{2}\\
    &= (\sqrt{5}+1) \cdot \frac{3-\sqrt{5}}{2}\\
    &= \frac{(\sqrt{5}+1)(3-\sqrt{5})}{2}\\
    &= \frac{3\sqrt{5} - 5 + 3 - \sqrt{5}}{2}\\
    &= \frac{2\sqrt{5} - 2}{2}\\
    &= \sqrt{5} - 1
    \end{align*}
\end{itemize}

\subsection*{B. Alternative Techniques}
\begin{itemize}[leftmargin=*]
    \item \textbf{Direct Coordinate Calculation}:
    \begin{itemize}
        \item Place pentagon with center at origin
        \item Vertices at $r_b(\cos(72k^\circ), \sin(72k^\circ))$ for $k=0,1,2,3,4$
        \item Calculate fold line equations
        \item Find intersection points forming small pentagon
        \item Calculate area directly
        \item More computational but avoids needing to memorize $\cos(36^\circ)$
    \end{itemize}
    
    \item \textbf{Area Subtraction Method}:
    \begin{itemize}
        \item Start with original pentagon area
        \item Subtract areas of folded triangular regions
        \item Add back overlapping regions (careful with double counting)
        \item Requires careful geometric analysis of the folded configuration
    \end{itemize}
\end{itemize}

\subsection*{C. Technique Selection Justification}
\begin{itemize}[leftmargin=*]
    \item \textbf{Why Similarity Method}: 
    \begin{itemize}
        \item Most elegant and direct
        \item Exploits the regularity and symmetry
        \item Requires only one key insight: $a_s = \frac{r_b}{2}$
        \item Uses standard pentagon property $\cos(36^\circ) = \frac{a}{r}$
    \end{itemize}
    \item \textbf{Time Efficiency}: 8-10 minutes if familiar with pentagon properties
\end{itemize}
\end{techniquesbox}

\begin{verificationbox}
\subsection*{A. Verification Level Selection}
\begin{itemize}[leftmargin=*]
    \item \textbf{Chosen Level}: Level 3 - Targeted Spot Check (20-30 seconds)
    \item \textbf{Justification}: 
    \begin{itemize}
        \item Late-band problem (P25)
        \item Geometric reasoning is sound once understood
        \item Can verify key algebraic steps
        \item Target confidence: 85-90\%
    \end{itemize}
    \item \textbf{Time Budget}: 25-30 seconds for verification
\end{itemize}

\subsection*{B. Verification Checklist}
\begin{itemize}[leftmargin=*]
    \item \textbf{Verify geometric setup}:
    \begin{itemize}
        \item Folding creates perpendicular bisector \checkmark
        \item Midpoint relationship: $a_s = \frac{r_b}{2}$ \checkmark
        \item Symmetry ensures regular pentagon \checkmark
    \end{itemize}
    \item \textbf{Verify scaling calculation}:
    \begin{itemize}
        \item $r_s = \frac{2r_b}{1+\sqrt{5}}$ \checkmark
        \item Rationalized: $\frac{r_s}{r_b} = \frac{\sqrt{5}-1}{2}$ \checkmark
        \item This equals $\frac{1}{\phi}$ (reciprocal of golden ratio) \checkmark
    \end{itemize}
    \item \textbf{Verify area ratio}:
    \begin{itemize}
        \item $\left(\frac{\sqrt{5}-1}{2}\right)^2 = \frac{6-2\sqrt{5}}{4} = \frac{3-\sqrt{5}}{2}$ \checkmark
    \end{itemize}
    \item \textbf{Verify final multiplication}:
    \begin{itemize}
        \item $(\sqrt{5}+1) \cdot \frac{3-\sqrt{5}}{2}$
        \item Expand: $\frac{3\sqrt{5} - 5 + 3 - \sqrt{5}}{2} = \frac{2\sqrt{5}-2}{2} = \sqrt{5}-1$ \checkmark
    \end{itemize}
    \item \textbf{Check reasonableness}:
    \begin{itemize}
        \item Original area: $\sqrt{5}+1 \approx 2.236 + 1 = 3.236$
        \item New area: $\sqrt{5}-1 \approx 2.236 - 1 = 1.236$
        \item Ratio: $\frac{1.236}{3.236} \approx 0.382 \approx \frac{1}{\phi}$ \checkmark
        \item Makes sense: smaller pentagon is significantly reduced
    \end{itemize}
    \item \textbf{Check answer}: Matches choice (B) \checkmark
\end{itemize}

\subsection*{C. Alternative Verification Methods}
\begin{itemize}[leftmargin=*]
    \item \textbf{Method 1}: Verify $\cos(36^\circ) = \frac{1+\sqrt{5}}{4}$
    \begin{itemize}
        \item Can derive from $\cos(72^\circ) = \frac{\sqrt{5}-1}{4}$ using double angle
        \item Or from golden ratio triangle with angles $36^\circ-72^\circ-72^\circ$
        \item Confirms the pentagon property used
    \end{itemize}
    \item \textbf{Method 2}: Check dimensions
    \begin{itemize}
        \item If scaling is $\frac{\sqrt{5}-1}{2} \approx 0.618$
        \item Area scaling should be $(0.618)^2 \approx 0.382$
        \item Actual: $\frac{\sqrt{5}-1}{\sqrt{5}+1} = \frac{(\sqrt{5}-1)^2}{4} = \frac{6-2\sqrt{5}}{4} \approx 0.382$ \checkmark
    \end{itemize}
    \item \textbf{Method 3}: Eliminate other choices
    \begin{itemize}
        \item (A) $4-\sqrt{5} \approx 1.764$ - too large
        \item (C) $8-3\sqrt{5} \approx 1.292$ - close but not exact
        \item (D) $\frac{\sqrt{5}+1}{2} \approx 1.618$ - too large (this is $\phi$)
        \item (E) $\frac{2+\sqrt{5}}{3} \approx 1.412$ - doesn't match our calculation
        \item Only (B) $\sqrt{5}-1$ matches
    \end{itemize}
\end{itemize}
\end{verificationbox}

\begin{solutionbox}
\subsection*{A. Step-by-Step Solution}
\begin{enumerate}[leftmargin=*]
    \item \textbf{Visualize the Folding}: 
    \begin{itemize}
        \item Regular pentagon $ABCDE$ with center $O$ and circumradius $r_b$
        \item Each vertex is folded to meet at center $O$
        \item Fold lines are perpendicular bisectors of radii
        \item By symmetry, fold lines form a regular pentagon (the answer)
    \end{itemize}
    
    \item \textbf{Key Geometric Insight}: 
    \begin{itemize}
        \item When vertex $A$ folds to $O$, fold line passes through midpoint $M$ of $OA$
        \item This fold line is an edge of the small pentagon
        \item Distance from $O$ to fold line = apothem of small pentagon = $a_s$
        \item Since fold line $\perp OA$ and passes through $M$: $a_s = OM = \frac{r_b}{2}$
    \end{itemize}
    
    \item \textbf{Apply Pentagon Property}: 
    \begin{itemize}
        \item For regular pentagon: $\cos(36^\circ) = \frac{\text{apothem}}{\text{circumradius}}$
        \item Known value: $\cos(36^\circ) = \frac{1+\sqrt{5}}{4}$
        \item For small pentagon: $\frac{a_s}{r_s} = \frac{1+\sqrt{5}}{4}$
        \item Therefore: $r_s = \frac{4a_s}{1+\sqrt{5}} = \frac{4 \cdot \frac{r_b}{2}}{1+\sqrt{5}} = \frac{2r_b}{1+\sqrt{5}}$
    \end{itemize}
    
    \item \textbf{Calculate Scaling Ratio}: 
    \begin{align*}
    \frac{r_s}{r_b} &= \frac{2}{1+\sqrt{5}} \cdot \frac{1-\sqrt{5}}{1-\sqrt{5}}\\
    &= \frac{2(1-\sqrt{5})}{1-5} = \frac{2(1-\sqrt{5})}{-4} = \frac{\sqrt{5}-1}{2}
    \end{align*}
    
    \item \textbf{Calculate Area Ratio}: 
    \begin{align*}
    \frac{A_s}{A_b} &= \left(\frac{\sqrt{5}-1}{2}\right)^2 = \frac{6-2\sqrt{5}}{4} = \frac{3-\sqrt{5}}{2}
    \end{align*}
    
    \item \textbf{Find Small Pentagon Area}: 
    \begin{align*}
    A_s &= (\sqrt{5}+1) \cdot \frac{3-\sqrt{5}}{2}\\
    &= \frac{(\sqrt{5}+1)(3-\sqrt{5})}{2}\\
    &= \frac{3\sqrt{5} - 5 + 3 - \sqrt{5}}{2}\\
    &= \frac{2\sqrt{5} - 2}{2} = \sqrt{5} - 1
    \end{align*}
\end{enumerate}

\subsection*{B. Alternative Approaches}

\textbf{Method 2: Using Complementary Golden Ratio Identity}

Once we establish that scaling ratio is $\frac{\sqrt{5}-1}{2} = \frac{1}{\phi}$:
\begin{itemize}
    \item Area ratio = $\frac{1}{\phi^2}$
    \item Using $\phi^2 = \phi + 1$: $\frac{1}{\phi^2} = \frac{1}{\phi+1}$
    \item Multiply: $\frac{1}{\phi+1} \cdot \frac{\phi}{\phi} = \frac{\phi}{\phi^2+\phi} = \frac{\phi}{\phi(\phi+1)+\phi} = ...$
    \item Direct calculation easier: $A_s = \frac{\sqrt{5}+1}{\phi^2}$
    \item Since $\phi^2 = \frac{3+\sqrt{5}}{2}$: $A_s = \frac{\sqrt{5}+1}{\frac{3+\sqrt{5}}{2}} = \frac{2(\sqrt{5}+1)}{3+\sqrt{5}} = \sqrt{5}-1$
\end{itemize}

\end{solutionbox}

\begin{competitionbox}
\subsection*{A. Time Management}
\begin{itemize}[leftmargin=*]
    \item \textbf{Estimated Time}: 8-12 minutes total
    \begin{itemize}
        \item Problem understanding and visualization: 1-2 minutes
        \item Identify folding property and $a_s = \frac{r_b}{2}$: 2-3 minutes
        \item Apply pentagon properties and calculate ratio: 3-4 minutes
        \item Calculate final area: 2-3 minutes
        \item Verification: 30 seconds
    \end{itemize}
    \item \textbf{Time Checkpoints}: 
    \begin{itemize}
        \item At 3 minutes: Should understand that $a_s = \frac{r_b}{2}$
        \item At 6 minutes: Should have scaling ratio $\frac{\sqrt{5}-1}{2}$
        \item At 10 minutes: Should have final answer
    \end{itemize}
    \item \textbf{Priority Level}: Low (skip unless very strong in geometry)
    \begin{itemize}
        \item P25 is optional for most students
        \item Only attempt if golden ratio and pentagon properties are familiar
        \item Skip if behind on time or uncertain about geometry
        \item Getting P16-P24 correct is more important
    \end{itemize}
\end{itemize}

\subsection*{B. Risk Assessment}
\begin{itemize}[leftmargin=*]
    \item \textbf{Confidence Level}: 90\% after full analysis
    \begin{itemize}
        \item Geometric reasoning is sound
        \item Algebra is careful and verified
        \item Answer matches choice (B)
        \item Numerical check confirms reasonableness
    \end{itemize}
    \item \textbf{Abandonment Criteria}: 
    \begin{itemize}
        \item If after 3 minutes can't visualize the folding, skip
        \item If don't know/can't derive $\cos(36^\circ) = \frac{1+\sqrt{5}}{4}$, skip
        \item If algebra becomes messy and time exceeds 10 minutes, move on
        \item Return at end only if time permits
    \end{itemize}
    \item \textbf{Flag for Review}: 
    \begin{itemize}
        \item Yes, if uncertain about the folding geometry
        \item No, if confident in similarity argument
    \end{itemize}
\end{itemize}

\subsection*{C. Answer Format}
\begin{itemize}[leftmargin=*]
    \item Multiple choice: Select (B) $\sqrt{5}-1$
    \item Double-check: Area is positive, less than original, involves $\sqrt{5}$ \checkmark
    \item Bubble carefully on answer sheet
\end{itemize}
\end{competitionbox}

\begin{keyinsightbox}
\textbf{Key Insight}: The crux of this problem is recognizing that when vertices are folded to the center, the fold lines pass through the midpoints of the radii. This creates the relationship:

$$\text{(apothem of small pentagon)} = \frac{1}{2} \times \text{(circumradius of large pentagon)}$$

Combined with the golden ratio property of regular pentagons:

$$\cos(36^\circ) = \frac{\text{apothem}}{\text{circumradius}} = \frac{1+\sqrt{5}}{4}$$

This gives us the scaling ratio $\frac{r_s}{r_b} = \frac{\sqrt{5}-1}{2} = \frac{1}{\phi}$ (reciprocal of golden ratio).

The area ratio is then $\frac{1}{\phi^2}$, which when applied to the original area $\sqrt{5}+1$ yields $\sqrt{5}-1$.

The beautiful appearance of the golden ratio $\phi$ is no accident—it's intrinsic to the geometry of regular pentagons!
\end{keyinsightbox}

\begin{warningbox}
\textbf{Warning}: Common pitfalls:
\begin{enumerate}
    \item \textbf{Misidentifying the fold line}: The fold is the perpendicular bisector, not just any line through the midpoint
    \item \textbf{Confusing apothem and circumradius}: Make sure to use $a_s = \frac{r_b}{2}$, not $r_s = \frac{r_b}{2}$
    \item \textbf{Wrong pentagon angle}: Use $36^\circ$ (half of $72^\circ$ central angle), not $54^\circ$ or $72^\circ$
    \item \textbf{Algebraic errors with $\sqrt{5}$}: Be very careful when expanding $(\sqrt{5}+1)(3-\sqrt{5})$
    \item \textbf{Forgetting to square the ratio}: Area ratio = (linear ratio)$^2$, not just linear ratio
    \item \textbf{Rationalizing errors}: When computing $\frac{2}{1+\sqrt{5}}$, multiply by $\frac{1-\sqrt{5}}{1-\sqrt{5}}$, not $\frac{\sqrt{5}-1}{\sqrt{5}-1}$
\end{enumerate}
\end{warningbox}

\begin{tipbox}
\textbf{Pro Tip}: For regular pentagon problems:

\textbf{Key formulas to know or derive}:
\begin{itemize}
    \item $\cos(36^\circ) = \sin(54^\circ) = \frac{1+\sqrt{5}}{4} = \frac{\phi}{2}$
    \item $\cos(72^\circ) = \sin(18^\circ) = \frac{\sqrt{5}-1}{4}$
    \item Golden ratio: $\phi = \frac{1+\sqrt{5}}{2}$, $\phi^2 = \phi+1$, $\frac{1}{\phi} = \phi-1 = \frac{\sqrt{5}-1}{2}$
\end{itemize}

\textbf{How to derive $\cos(36^\circ)$ if forgotten}:
\begin{enumerate}
    \item Consider isosceles triangle with angles $36^\circ-72^\circ-72^\circ$ (appears in regular pentagon)
    \item Let equal sides have length 1, base has length $x$
    \item Bisect one base angle to create $36^\circ-36^\circ-108^\circ$ triangle
    \item By similarity: $\frac{1}{x} = \frac{x}{1-x}$, giving $x^2+x-1=0$
    \item Solve: $x = \frac{\sqrt{5}-1}{2}$ (positive root)
    \item Use law of cosines to find $\cos(36^\circ) = \frac{1+\sqrt{5}}{4}$
\end{enumerate}

\textbf{Paper folding insight}: Whenever a point is folded onto another point, the fold line is the perpendicular bisector of the segment connecting them. This is a powerful geometric principle applicable beyond this problem.

\textbf{Competition strategy}: P25 problems often involve special properties of regular polygons, golden ratios, or other elegant mathematical relationships. If you recognize these patterns, the problem becomes much more tractable. If not familiar with these properties, it's often wise to skip and ensure earlier problems are correct.
\end{tipbox}

\section*{Final Answer}

By establishing that the folding operation creates a perpendicular bisector relationship where $a_s = \frac{r_b}{2}$, and using the golden ratio property of regular pentagons $\cos(36^\circ) = \frac{1+\sqrt{5}}{4}$, we find:

\textbf{Scaling ratio}: $\frac{r_s}{r_b} = \frac{\sqrt{5}-1}{2} = \frac{1}{\phi}$

\textbf{Area ratio}: $\left(\frac{\sqrt{5}-1}{2}\right)^2 = \frac{3-\sqrt{5}}{2}$

\textbf{Small pentagon area}: 
$$A_s = (\sqrt{5}+1) \cdot \frac{3-\sqrt{5}}{2} = \sqrt{5}-1$$

$$\boxed{\textbf{(B) } \sqrt{5}-1}$$

\section*{Confidence Assessment}
\begin{itemize}[leftmargin=*]
    \item \textbf{Confidence Level}: 90\% 
    \begin{itemize}
        \item Geometric reasoning is rigorous: folding creates perpendicular bisector
        \item Key relationship $a_s = \frac{r_b}{2}$ follows from geometry
        \item Pentagon property $\cos(36^\circ) = \frac{1+\sqrt{5}}{4}$ is standard
        \item Algebraic calculations verified at each step
        \item Answer matches choice (B) and passes reasonableness check
        \item Golden ratio relationship $\frac{1}{\phi}$ is elegant and expected for pentagons
    \end{itemize}
    \item \textbf{Verification Applied}: Level 3 - Targeted Spot Check
    \begin{itemize}
        \item Primary method: Similarity with folding geometry
        \item Verified: Each algebraic manipulation step
        \item Cross-check: Numerical approximation confirms ratio
        \item Consistency: Answer has expected form with $\sqrt{5}$
    \end{itemize}
    \item \textbf{Flag for Review}: No
    \begin{itemize}
        \item High confidence after systematic analysis
        \item Geometry and algebra both solid
        \item Multiple verification methods confirm answer
    \end{itemize}
    \item \textbf{Time Spent}: 
    \begin{itemize}
        \item Estimated: 8-10 minutes for complete solution
        \item Appropriate for P25 with geometric insight
        \item Golden ratio knowledge saves time
    \end{itemize}
\end{itemize}

\end{document}

